% Chương này đã được khai báo trong main.tex.

\section{Mô hình hóa bài toán vận tải}
    % section: Mô hình Bài toán Vận tải

\begin{table}[H]
    \centering
    \begin{tabular}{|c|c|c|c|}
        \hline
        & \textbf{Trạm thu $B_1$} & \textbf{Trạm thu $B_2$} & \textbf{Trạm thu $B_3$} \\
        \hline
        \textbf{Trạm phát $A_1$} & $c_{11}$ & $c_{12}$ & $c_{13}$ \\
        \hline
        \textbf{Trạm phát $A_2$} & $c_{21}$ & $c_{22}$ & $c_{23}$ \\
        \hline
        \textbf{Trạm phát $A_3$} & $c_{31}$ & $c_{32}$ & $c_{33}$ \\
        \hline
    \end{tabular}
    \caption{Cước phí vận chuyển từ các trạm phát đến các trạm thu}
\end{table}

\begin{figure}[H]
    \centering
    \resizebox{0.95\textwidth}{!}{%
    \begin{tikzpicture}[
        >={Stealth[scale=1.2]},
        cell/.style={rectangle, draw=black, minimum width=3.5cm, minimum height=1.5cm, align=center, anchor=center},
        ampersand replacement=\&
    ]
        % 1. Vẽ ma trận 3x4
        \matrix (M) [matrix of nodes, nodes={cell}, column sep=-\pgflinewidth, row sep=-\pgflinewidth, nodes in empty cells]
        {
            |[alias=diag]| \& |[alias=b1]| $T_1 : t_1$ \& |[alias=b2]| $T_2 : t_2$ \& |[alias=b3]| $T_3 : t_3$ \\
            |[alias=a1]| $P_1 : p_1$ \& |[alias=cell11]| $c_{11} \mid x_{11}$ \& $c_{12} \mid x_{12}$ \& $c_{13} \mid x_{13}$ \\
            |[alias=a2]| $P_2 : p_2$ \& $c_{21} \mid x_{21}$ \& $c_{22} \mid x_{22}$ \& $c_{23} \mid x_{23}$ \\
        };

        % 2. Xử lý ô 1x1 (Đường chéo và Label)
        % Đường chéo
        \draw[black!70, thick] (diag.north west) -- (diag.south east);

        % Trạm thu (trên phải)
        \node[
            anchor=north east,
            xshift=-3pt,
            yshift=-3pt,
            font=\footnotesize,
            fill=white,
            inner sep=1pt
        ] at (diag.north east)
        {Trạm thu ($T_j$)};

        % Trạm phát (dưới trái)
        \node[
            anchor=south west,
            xshift=3pt,
            yshift=3pt,
            font=\footnotesize,
            fill=white,
            inner sep=1pt
        ] at (diag.south west)
        {Trạm phát ($P_i$)};

        % 3. Khoanh vùng 1x2 -> 1x4 (Lượng hàng trạm thu)
        \node[draw=blue, thick, rounded corners, inner sep=-4pt, fit=(b1) (b3)] (groupT) {};
        \draw[<-, blue, thick] (groupT.north) -- ++(0,1) node[above, font=\bfseries] {Nhu cầu tại trạm thu $T_j$ ($T_j$)};

        % 4. Khoanh vùng 2x1 -> 3x1 (Lượng hàng trạm phát)
        \node[draw=orange, thick, rounded corners, inner sep=-4pt, fit=(a1) (a2)] (groupP) {};
        \draw[<-, orange, thick] (groupP.south) -- ++(0,-1.2) node[below, align=center, font=\bfseries] {Lượng phát tại trạm $P_i$ ($p_i$)};

        % 5. Mũi tên giải thích cho c_ij và x_ij (Lấy ô cell11 làm chuẩn)
        % Tạo node ảo lùi về trái/phải để mũi tên chỉ đúng vào c_11 và x_11
        \node at (cell11.center) [xshift=-0.6cm] (cij) {};
        \node at (cell11.center) [xshift=0.6cm] (xij) {};

        % Mũi tên cho c_ij (Cước phí)
        \draw[<-, red, ultra thick] (cij) -- ++(2,-2.6) node[below, align=center, text=red, font=\small] {Chi phí vận chuyển một đơn vị \\ từ trạm $P_i$ đến trạm $T_j$};
        
        % Mũi tên cho x_ij (Lượng hàng)
        % 
        \draw[<-, purple, ultra thick] (xij) -- ++(5,-3.5) node[below, align=center, text=purple, font=\small] {Lượng hàng vận chuyển thực tế \\ từ trạm $P_i$ đến trạm $T_j$};

    \end{tikzpicture}
    }
    \caption{Mô hình hóa Bảng phân phát bài toán vận tải}
    \label{fig:mohinh-phanphat-vantai}
\end{figure}

\begin{figure}[H]
    \centering
    \resizebox{0.95\textwidth}{!}{%
    \begin{tikzpicture}[
        >={Stealth[scale=1.2]},
        cell/.style={rectangle, draw=black, minimum width=3.5cm, minimum height=1.5cm, align=center, anchor=center},
        ampersand replacement=\&
    ]
        % 1. Vẽ ma trận 3x4
        \matrix (M) [matrix of nodes, nodes={cell}, column sep=-\pgflinewidth, row sep=-\pgflinewidth, nodes in empty cells]
        {
            |[alias=diag]| \& |[alias=b1]| $300 : t_1$ \& |[alias=b2]| $500 : t_2$ \& |[alias=b3]| $100 : t_3$ \\
            |[alias=a1]| $700 : p_1$ \& |[alias=cell11]| $5 \mid x_{11}$ \& $6 \mid x_{12}$ \& $4 \mid x_{13}$ \\
            |[alias=a2]| $200 : p_2$ \& $8 \mid x_{21}$ \& $9 \mid x_{22}$ \& $5 \mid x_{23}$ \\
        };

        % 2. Xử lý ô 1x1 (Đường chéo và Label)
        % Đường chéo
        \draw[black!70, thick] (diag.north west) -- (diag.south east);

        % Trạm thu (trên phải)
        \node[
            anchor=north east,
            xshift=-3pt,
            yshift=-3pt,
            font=\footnotesize,
            fill=white,
            inner sep=1pt
        ] at (diag.north east)
        {Trạm thu ($T_j$)};

        % Trạm phát (dưới trái)
        \node[
            anchor=south west,
            xshift=3pt,
            yshift=3pt,
            font=\footnotesize,
            fill=white,
            inner sep=1pt
        ] at (diag.south west)
        {Trạm phát ($P_i$)};

    \end{tikzpicture}
    }
    \caption{Ví dụ Về Mô hình hóa Bảng phân phát bài toán vận tải}
    \label{fig:mohinh-phanphat-vantai}
\end{figure}


Nhận xét: \\
$$\sum \text{phát} = 700 + 200 = 900$$
$$\sum \text{thu} = 300 + 500 + 100 = 900$$
$$=> \sum \text{phát} = \sum \text{thu}$$
Trong mô hình bài toán vận tải, chúng ta có một tập hợp các trạm phát (các điểm cung cấp hàng hóa) và một tập hợp các trạm thu (các điểm tiêu thụ hàng hóa). Mỗi trạm phát có một lượng hàng hóa cụ thể cần được vận chuyển đến các trạm thu, và mỗi trạm thu có một nhu cầu cụ thể cần được đáp ứng. Cước phí vận chuyển từ mỗi trạm phát đến mỗi trạm thu được biểu diễn bằng một ma trận chi phí $c_{ij}$, trong đó $i$ là chỉ số của trạm phát và $j$ là chỉ số của trạm thu. Mục tiêu của bài toán là tìm ra phương án vận chuyển tối ưu sao cho tổng chi phí vận chuyển được tối thiểu hóa, đồng thời đảm bảo rằng tất cả các nhu cầu của các trạm thu được đáp ứng và không vượt quá khả năng cung cấp của các trạm phát.

\section{Các khái niệm}
\subsection{Các khái niệm về cấu trúc mô hình bài toán}
\begin{itemize}
    \item \textbf{Trạm phát (Supply Point):} Là các điểm cung cấp hàng hóa, thường được ký hiệu là $A_i$ với $i = \overline{1,m}$, mỗi trạm phát có một lượng hàng hóa cụ thể $a_i$ (trữ lượng).
    \item \textbf{Trạm thu (Demand Point):} Là các điểm tiêu thụ hàng hóa, thường được ký hiệu là $B_j$ với $j = \overline{1,n}$, mỗi trạm thu có một nhu cầu cụ thể $b_j$.
    \item \textbf{Cước phí vận chuyển (Transportation Cost):} Là chi phí hoặc hiệu suất liên quan đến việc vận chuyển một đơn vị hàng hóa từ trạm phát $A_i$ đến trạm thu $B_j$, được ký hiệu là $c_{ij}$.
    \item \textbf{Biến quyết định (Decision Variable):} Là các biến $x_{ij}$ đại diện cho lượng hàng hóa được vận chuyển từ trạm phát $A_i$ đến trạm thu $B_j$, cần được xác định để tối ưu hóa mục tiêu của bài toán.
    \item \textbf{Hàm mục tiêu (Objective Function):} Là hàm cần được tối ưu hóa (tối thiểu hóa tổng chi phí hoặc tối đa hóa hiệu suất), thường được biểu diễn dưới dạng $Z = \sum_{i=1}^{m} \sum_{j=1}^{n} c_{ij} x_{ij}$.
\end{itemize}
\subsection{Các khái niệm về Phương án và Ô phân bổ}
\begin{itemize}
    \item \textbf{Phương án (Solution):} Là một tập hợp các giá trị $x_{ij}$ đại diện cho lượng hàng hóa được vận chuyển từ trạm phát $A_i$ đến trạm thu $B_j$. Một phương án phải thỏa mãn tất cả các ràng buộc về cung cấp và nhu cầu.
    \item \textbf{Ô phân bổ (Basic Variable):} Là các biến $x_{ij}$ có giá trị dương trong một phương án cơ bản. Số lượng ô phân bổ tối đa là $m + n - 1$ để đảm bảo tính khả thi của phương án.
    \item \textbf{Ô không phân bổ (Non-basic Variable):} Là các biến $x_{ij}$ có giá trị bằng 0 trong một phương án cơ bản. Các ô không phân bổ không đóng góp vào tổng chi phí hoặc hiệu suất của phương án.
\end{itemize}


\subsection{Phân loại Phương án Cơ bản}
\begin{itemize}
    \item \textbf{Phương án cơ bản (PACB)}: Là phương án mà các ô chọn cấu thành nên nó không tạo ra bất kỳ chu trình nào.
    \item \textbf{PACB không suy biến}: Xảy ra khi số lượng ô chọn trong phương án đạt đúng số lượng $m + n - 1$ (với $m$ là số trạm phát và $n$ là số trạm thu), đảm bảo rằng không có ô nào bị suy biến (tức là không có ô nào có giá trị bằng 0).
    \item \textbf{PACB suy biến}: Xảy ra khi số lượng ô chọn trong phương án nhỏ hơn $m + n - 1$, dẫn đến việc có ít nhất một ô có giá trị bằng 0.
\end{itemize}

\section{Các phương pháp giải bài toán vận tải}
Quy tắc phân phối chung: Phân phối hàng hóa từ trạm phát đến trạm thu sao cho tổng chi phí vận chuyển được tối thiểu hóa, đồng thời đảm bảo rằng tất cả các nhu cầu của các trạm thu được đáp ứng và không vượt quá khả năng cung cấp của các trạm phát. \\
Tổng quát bằng công thức toán học:
\begin{formula}
    x_{ij} = \min(p_i^{(k)}, t_j^{(k)})
\end{formula}
Trong đó: 
\begin{itemize}
    \item $x_{ij}$: Lượng hàng hóa được phân bổ từ trạm phát $p_i$ đến trạm thu $t_j$.
    \item $p_i^{(k)}$: Lượng hàng hóa còn lại tại trạm phát $p_i$ sau khi đã phân bổ trong bước thứ $k$.
    \item $t_j^{(k)}$: Nhu cầu còn lại tại trạm thu $t_j$ sau khi đã phân bổ trong bước thứ $k$.
\end{itemize}

Cập nhật lại Cung và Cầu cho bước $k+1$:
$$ p_i^{(k+1)} = p_i^{(k)} - x_{ij} $$
$$ t_j^{(k+1)} = t_j^{(k)} - x_{ij} $$
Điều kiện gạch bỏ: 
\begin{itemize}
    \item Nếu $p_i^{(k+1)} = 0$, gạch bỏ trạm phát $p_i$ khỏi bảng phân phối. 
    \item Nếu $t_j^{(k+1)} = 0$, gạch bỏ trạm thu $t_j$ khỏi bảng phân phối.
\end{itemize}

    \subsection{Phương pháp chi phí thấp nhất (Least Cost Method)}

Phương pháp chi phí thấp nhất là một kỹ thuật giải quyết bài toán vận tải bằng cách ưu tiên phân bổ hàng hóa từ trạm phát đến trạm thu dựa trên cước phí vận chuyển thấp nhất. Dưới đây là các bước thực hiện phương pháp này:

\begin{enumerate}
    \item \textbf{Xác định chi phí thấp nhất:} Tìm ô có cước phí vận chuyển $c_{ij}$ thấp nhất trong phần bảng chưa bị gạch bỏ. Nếu có nhiều ô cùng chi phí thấp nhất, ưu tiên chọn ô ở góc Tây - Bắc (trên cùng bên trái).
    
    \item \textbf{Phân bổ hàng hóa:} Phân bổ lượng hàng hóa tối đa có thể vào ô vừa chọn theo công thức: $x_{ij} = \min(a_i, b_j)$ (với $a_i$ là khả năng cung cấp còn lại của trạm phát $P_i$, và $b_j$ là nhu cầu còn lại của trạm thu $T_j$).
    
    \item \textbf{Cập nhật và Gạch bỏ:} 
    \begin{itemize}
        \item Cập nhật lại lượng cung: $a_i \leftarrow a_i - x_{ij}$.
        \item Cập nhật lại lượng cầu: $b_j \leftarrow b_j - x_{ij}$.
        \item Trạm nào (phát hoặc thu) giảm về 0 trước thì gạch bỏ dòng hoặc cột tương ứng khỏi bảng để không xét trong các bước tiếp theo.
    \end{itemize}
    
    \item \textbf{Lặp lại:} Lặp lại quá trình trên cho đến khi tất cả nhu cầu của các trạm thu được đáp ứng (tổng cung và tổng cầu đều về 0).
\end{enumerate}

\textbf{Lưu ý:} Phương pháp chi phí thấp nhất không đảm bảo rằng phương án cơ bản xuất phát thu được sẽ là phương án tối ưu.

\textbf{Tổng quát Phương pháp}:
\begin{itemize}
    \item Phân phối lượng hàng tôi đa có thể (Dựa theo Công thức ) vào ô có chi phí thấp nhất ($\min(c_{ij})$) trước tiên.
    \item Cập nhật lại cung và cầu sau mỗi bước phân bổ, và gạch bỏ các trạm đã hoàn thành (cung hoặc cầu về 0) để thu hẹp phạm vi xét tiếp theo.
    \item Trường hợp $> 1$ ô $c_{ij} = \min(c_{ij})$ -> Chọn bất kỳ.
\end{itemize}

\begin{figure}[H]
    \centering
    \resizebox{0.95\textwidth}{!}{%
    \begin{tikzpicture}[
        >={Stealth[scale=1.2]},
        line cap=rect,
        line join=miter,
        cell/.style={
            rectangle,
            draw=black,
            line width=0.5pt,
            minimum width=3.5cm,
            minimum height=1.5cm,
            align=center,
            anchor=center,
            inner sep=2pt
        },
        ampersand replacement=\&
    ]
        % 1. Vẽ ma trận 3x4 (Đã bổ sung đầy đủ alias cho các ô)
        \matrix (M) [
            matrix of nodes,
            nodes={cell},
            nodes in empty cells,
            column sep=-0.5pt,
            row sep=-0.5pt
        ]
        {
            |[alias=diag]| \& |[alias=b1]| $\cancel{300} \rightarrow 0 : T_1$ \& |[alias=b2]| $\cancel{500} \rightarrow \cancel{200} \rightarrow 0 : T_2$ \& |[alias=b3]| $\cancel{100} \rightarrow 0 : T_3$ \\
            |[alias=a1]| $\cancel{700} \rightarrow \cancel{600} \rightarrow \cancel{300} \rightarrow 0: P_1$ \& |[alias=cell11]| $5 \mid 300$ \& |[alias=cell12]| $6 \mid 300$ \& |[alias=cell13]| $4 \mid 100$ \\
            |[alias=a2]| $\cancel{200} \rightarrow 0 : P_2$ \& |[alias=cell21]| $8 \mid 0$ \& |[alias=cell22]| $9 \mid 200$ \& |[alias=cell23]| $5 \mid 0$ \\
        };

        % Ký hiệu x_ij ở góc dưới-trái từng ô
        \node[anchor=south west, xshift=2pt, yshift=2pt, font=\scriptsize] at (cell11.south west) {$x_{11}$};
        \node[anchor=south west, xshift=2pt, yshift=2pt, font=\scriptsize] at (cell12.south west) {$x_{12}$};
        \node[anchor=south west, xshift=2pt, yshift=2pt, font=\scriptsize] at (cell13.south west) {$x_{13}$};
        \node[anchor=south west, xshift=2pt, yshift=2pt, font=\scriptsize] at (cell21.south west) {$x_{21}$};
        \node[anchor=south west, xshift=2pt, yshift=2pt, font=\scriptsize] at (cell22.south west) {$x_{22}$};
        \node[anchor=south west, xshift=2pt, yshift=2pt, font=\scriptsize] at (cell23.south west) {$x_{23}$};

        % 2. Xử lý ô 1x1 (Đường chéo và Label)
        % Đường chéo
        \draw[black!70, thick] (diag.north west) -- (diag.south east);

        % Trạm thu (trên phải)
        \node[
            anchor=north east,
            xshift=-3pt,
            yshift=-3pt,
            font=\footnotesize,
            fill=white,
            inner sep=1pt
        ] at (diag.north east)
        {Trạm thu ($T_j$)};

        % Trạm phát (dưới trái)
        \node[
            anchor=south west,
            xshift=3pt,
            yshift=3pt,
            font=\footnotesize,
            fill=white,
            inner sep=1pt
        ] at (diag.south west)
        {Trạm phát ($P_i$)};

        % (Tùy chọn) Khoanh đỏ ô bắt đầu bằng TikZ tiêu chuẩn thay vì \koval
        \draw[red, thick, rounded corners] ([xshift=5pt,yshift=-5pt]cell13.north west) rectangle ([xshift=-5pt,yshift=5pt]cell13.south east);

    \end{tikzpicture}
    }
    \caption{Ví dụ Phương pháp chi phí thấp nhất trong bài toán vận tải}
    \label{fig:least-cost-method-transport}
\end{figure}

\textbf{Diễn giải chi tiết các bước điền:}
\begin{enumerate}
    \item Tìm $\min c_{ij}$ trên toàn bảng: $\min c_{ij} = c_{13} = 4$.
    \begin{itemize}
        \item Phân bổ: $x_{13} = \min(b_3^{(1)}, a_1^{(1)}) = \min(100, 700) = 100$.
        \item Cập nhật: $a_1^{(2)} = 700 - 100 = 600$; $b_3^{(2)} = 100 - 100 = 0$.
        \item Do $b_3 = 0$, ta \textbf{gạch bỏ cột $T_3$} khỏi các bước xét tiếp theo.
    \end{itemize}
    
    \item Tìm $\min c_{ij}$ trong các ô còn lại (cột 1 và cột 2): $\min c_{ij} = c_{11} = 5$.
    \begin{itemize}
        \item Phân bổ: $x_{11} = \min(b_1^{(1)}, a_1^{(2)}) = \min(300, 600) = 300$.
        \item Cập nhật: $a_1^{(3)} = 600 - 300 = 300$; $b_1^{(2)} = 300 - 300 = 0$.
        \item Do $b_1 = 0$, ta \textbf{gạch bỏ cột $T_1$}.
    \end{itemize}
    
    \item Tìm $\min c_{ij}$ trong các ô còn lại (chỉ còn cột 2): $\min c_{ij} = c_{12} = 6$.
    \begin{itemize}
        \item Phân bổ: $x_{12} = \min(b_2^{(1)}, a_1^{(3)}) = \min(500, 300) = 300$.
        \item Cập nhật: $a_1^{(4)} = 300 - 300 = 0$; $b_2^{(2)} = 500 - 300 = 200$.
        \item Do $a_1 = 0$, ta \textbf{gạch bỏ dòng $P_1$}.
    \end{itemize}
    
    \item Ô duy nhất còn lại chưa bị gạch là $x_{22}$ với $c_{22} = 9$.
    \begin{itemize}
        \item Phân bổ: $x_{22} = \min(b_2^{(2)}, a_2^{(1)}) = \min(200, 200) = 200$.
        \item Cập nhật: $a_2^{(2)} = 0$; $b_2^{(3)} = 0$.
        \item Tất cả cung và cầu đều đã về 0. Quá trình kết thúc.
    \end{itemize}
\end{enumerate}

\textbf{Kết luận:} Phương án cơ bản xuất phát thu được là:
\[ X^0 = \begin{pmatrix} 300 & 300 & 100 \\ 0 & 200 & 0 \end{pmatrix} \]
$X^0$ có 4 ô chọn = $m + n - 1 = 2 + 3 - 1 = 4$ \\
$=> X^0$ là phương án cơ bản không suy biến \\
Với tổng cước phí vận chuyển ban đầu là: 
$Z = (300 \times 5) + (300 \times 6) + (100 \times 4) + (200 \times 9) = 1500 + 1800 + 400 + 1800 = 5500$.
    \subsection{Phương pháp góc Tây Bắc (Northwest Corner Method)}

Phương pháp góc Tây Bắc là một kỹ thuật giải quyết bài toán vận tải bằng cách ưu tiên phân bổ hàng hóa vào ô nằm ở vị trí trên cùng bên trái (góc Tây Bắc) của bảng phân phát, hoàn toàn không quan tâm đến cước phí vận chuyển. Dưới đây là các bước thực hiện phương pháp này:

\begin{enumerate}
    \item \textbf{Xác định ô góc Tây Bắc:} Luôn chọn ô nằm ở vị trí trên cùng bên trái của phần bảng phân phát chưa bị gạch bỏ (ô có chỉ số dòng và cột nhỏ nhất còn lại).
    
    \item \textbf{Phân bổ hàng hóa:} Phân bổ lượng hàng hóa tối đa có thể vào ô vừa chọn theo công thức: $x_{ij} = \min(a_i, b_j)$ (với $a_i$ là khả năng cung cấp còn lại của trạm phát $P_i$, và $b_j$ là nhu cầu còn lại của trạm thu $T_j$).
    
    \item \textbf{Cập nhật và Gạch bỏ:} 
    \begin{itemize}
        \item Cập nhật lại lượng cung: $a_i \leftarrow a_i - x_{ij}$.
        \item Cập nhật lại lượng cầu: $b_j \leftarrow b_j - x_{ij}$.
        \item Trạm nào (phát hoặc thu) giảm về 0 trước thì gạch bỏ dòng hoặc cột tương ứng khỏi bảng. Nếu cả cung và cầu cùng về 0 đồng thời, ưu tiên gạch một chiều (dòng hoặc cột), chiều còn lại giữ nguyên giá trị 0 để xét tiếp.
    \end{itemize}
    
    \item \textbf{Lặp lại:} Tiếp tục lặp lại quá trình trên (luôn chọn góc Tây Bắc của phần bảng còn lại) cho đến khi tất cả nhu cầu và cung cấp đều về 0.
\end{enumerate}

\textbf{Lưu ý:} Phương pháp góc Tây Bắc rất dễ thực hiện và có tính hệ thống cao. Tuy nhiên, do hoàn toàn bỏ qua yếu tố chi phí, phương án xuất phát thu được thường có tổng cước phí khá cao và tốn nhiều bước lặp hơn để tối ưu bằng phương pháp Thế vị.

\textbf{Tổng quát Phương pháp}:
\begin{itemize}
    \item Luôn ưu tiên ô $(i, j)$ có chỉ số nhỏ nhất.
    \item Điền lượng hàng phân phối $x_{ij} = \min(a_i, b_j)$.
    \item Cập nhật lại cung và cầu, gạch bỏ các trạm đã hoàn thành và "tịnh tiến" sang cột bên phải hoặc xuống dòng bên dưới.
\end{itemize}

\begin{figure}[H]
    \centering
    \resizebox{0.95\textwidth}{!}{%
    \begin{tikzpicture}[
        >={Stealth[scale=1.2]},
        line cap=rect,
        line join=miter,
        cell/.style={
            rectangle,
            draw=black,
            line width=0.5pt,
            minimum width=3.5cm,
            minimum height=1.5cm,
            align=center,
            anchor=center,
            inner sep=2pt
        },
        ampersand replacement=\&
    ]
        % 1. Vẽ ma trận 3x4
        \matrix (M) [
            matrix of nodes,
            nodes={cell},
            nodes in empty cells,
            column sep=-0.5pt,
            row sep=-0.5pt
        ]
        {
            |[alias=diag]| \& |[alias=b1]| $\cancel{300} \rightarrow 0 : T_1$ \& |[alias=b2]| $\cancel{500} \rightarrow \cancel{100} \rightarrow 0 : T_2$ \& |[alias=b3]| $\cancel{100} \rightarrow 0 : T_3$ \\
            |[alias=a1]| $\cancel{700} \rightarrow \cancel{400} \rightarrow 0: P_1$ \& |[alias=cell11]| $5 \mid 300$ \& |[alias=cell12]| $6 \mid 400$ \& |[alias=cell13]| $4 \mid 0$ \\
            |[alias=a2]| $\cancel{200} \rightarrow \cancel{100} \rightarrow 0 : P_2$ \& |[alias=cell21]| $8 \mid 0$ \& |[alias=cell22]| $9 \mid 100$ \& |[alias=cell23]| $5 \mid 100$ \\
        };

        % Ký hiệu x_ij ở góc dưới-trái từng ô
        \node[anchor=south west, xshift=2pt, yshift=2pt, font=\scriptsize] at (cell11.south west) {$x_{11}$};
        \node[anchor=south west, xshift=2pt, yshift=2pt, font=\scriptsize] at (cell12.south west) {$x_{12}$};
        \node[anchor=south west, xshift=2pt, yshift=2pt, font=\scriptsize] at (cell13.south west) {$x_{13}$};
        \node[anchor=south west, xshift=2pt, yshift=2pt, font=\scriptsize] at (cell21.south west) {$x_{21}$};
        \node[anchor=south west, xshift=2pt, yshift=2pt, font=\scriptsize] at (cell22.south west) {$x_{22}$};
        \node[anchor=south west, xshift=2pt, yshift=2pt, font=\scriptsize] at (cell23.south west) {$x_{23}$};

        % 2. Xử lý ô 1x1 (Đường chéo và Label)
        \draw[black!70, thick] (diag.north west) -- (diag.south east);
        \node[anchor=north east, xshift=-3pt, yshift=-3pt, font=\footnotesize, fill=white, inner sep=1pt] at (diag.north east) {Trạm thu ($T_j$)};
        \node[anchor=south west, xshift=3pt, yshift=3pt, font=\footnotesize, fill=white, inner sep=1pt] at (diag.south west) {Trạm phát ($P_i$)};

        % Khoanh đỏ ô xuất phát đầu tiên x_11 để nhấn mạnh
        \draw[red, thick, rounded corners] ([xshift=5pt,yshift=-5pt]cell11.north west) rectangle ([xshift=-5pt,yshift=5pt]cell11.south east);

    \end{tikzpicture}
    }
    \caption{Ví dụ Phương pháp góc Tây Bắc trong bài toán vận tải}
    \label{fig:nw-corner-method-transport}
\end{figure}

\textbf{Diễn giải chi tiết các bước điền:}
\begin{enumerate}
    \item Bắt đầu từ ô góc Tây Bắc của toàn bảng là $x_{11}$.
    \begin{itemize}
        \item Phân bổ: $x_{11} = \min(b_1^{(1)}, a_1^{(1)}) = \min(300, 700) = 300$.
        \item Cập nhật: $a_1^{(2)} = 700 - 300 = 400$; $b_1^{(2)} = 300 - 300 = 0$.
        \item Do $b_1 = 0$, ta \textbf{gạch bỏ cột $T_1$} khỏi các bước xét tiếp theo.
    \end{itemize}
    
    \item Ô góc Tây Bắc của phần bảng còn lại (từ cột 2) là $x_{12}$.
    \begin{itemize}
        \item Phân bổ: $x_{12} = \min(b_2^{(1)}, a_1^{(2)}) = \min(500, 400) = 400$.
        \item Cập nhật: $a_1^{(3)} = 400 - 400 = 0$; $b_2^{(2)} = 500 - 400 = 100$.
        \item Do $a_1 = 0$, ta \textbf{gạch bỏ dòng $P_1$}.
    \end{itemize}
    
    \item Ô góc Tây Bắc của phần bảng còn lại (chỉ còn dòng 2, cột 2 và 3) là $x_{22}$.
    \begin{itemize}
        \item Phân bổ: $x_{22} = \min(b_2^{(2)}, a_2^{(1)}) = \min(100, 200) = 100$.
        \item Cập nhật: $a_2^{(2)} = 200 - 100 = 100$; $b_2^{(3)} = 100 - 100 = 0$.
        \item Do $b_2 = 0$, ta \textbf{gạch bỏ cột $T_2$}.
    \end{itemize}
    
    \item Ô duy nhất còn lại ở góc Tây Bắc (và cũng là ô cuối cùng) là $x_{23}$.
    \begin{itemize}
        \item Phân bổ: $x_{23} = \min(b_3^{(1)}, a_2^{(2)}) = \min(100, 100) = 100$.
        \item Cập nhật: $a_2^{(3)} = 0$; $b_3^{(2)} = 0$.
        \item Tất cả cung và cầu đều đã về 0. Quá trình kết thúc.
    \end{itemize}
\end{enumerate}

\textbf{Kết luận:} Phương án cơ bản xuất phát thu được là:
\[ X^0 = \begin{pmatrix} 300 & 400 & 0 \\ 0 & 100 & 100 \end{pmatrix} \]
$X^0$ có 4 ô chọn = $m + n - 1 = 2 + 3 - 1 = 4$ \\
$\Rightarrow X^0$ là phương án cơ bản không suy biến. \\
Với tổng cước phí vận chuyển ban đầu là: 
$Z = (300 \times 5) + (400 \times 6) + (100 \times 9) + (100 \times 5) = 1500 + 2400 + 900 + 500 = 5300$.

\textit{Nhận xét:} Trong trường hợp bài toán cụ thể này, phương pháp góc Tây Bắc ngẫu nhiên cho ra kết quả ($Z = 5300$) tốt hơn so với phương pháp chi phí thấp nhất ($Z = 5500$) do tránh được việc phải dồn một lượng lớn hàng (200 đơn vị) vào ô có cước phí cực cao ($c_{22}=9$). Tuy nhiên, đây chỉ là một trường hợp cụ thể và không phải lúc nào phương pháp góc Tây Bắc cũng cho ra phương án tốt hơn. Do đó, việc lựa chọn phương pháp khởi tạo phù hợp với đặc điểm của bài toán là rất quan trọng để đạt được hiệu quả tối ưu nhanh chóng.
    \subsection{Phương pháp xấp xỉ Vogel (Vogel's Approximation Method - VAM / Fogels)}

Phương pháp Vogel (hay Fogels) là một kỹ thuật tìm phương án xuất phát tối ưu hơn bằng cách đánh giá chi phí cơ hội (độ chênh lệch cước phí) nếu không chọn ô có cước phí thấp nhất. Các bước thực hiện như sau:

\begin{enumerate}
    \item \textbf{Tính độ chênh lệch (Penalty - $\Delta$):} Trên mỗi dòng và mỗi cột chưa bị gạch bỏ, tìm hai ô có cước phí thấp nhất ($c_{min1}$) và thấp thứ hai ($c_{min2}$). Tính độ chênh lệch: $\Delta = c_{min2} - c_{min1}$ và ghi ra rìa ngoài của bảng. (Nếu chỉ còn 1 ô chưa gạch trên dòng/cột, $\Delta$ chính là cước phí của ô đó).
    
    \item \textbf{Chọn dòng/cột ưu tiên:} Chọn dòng hoặc cột có độ chênh lệch $\Delta$ lớn nhất (khoanh tròn lại). Nếu có nhiều $\Delta$ bằng nhau, ta có thể chọn tùy ý (thường ưu tiên dòng/cột chứa ô có cước phí $c_{min1}$ nhỏ hơn để tối ưu chi phí).
    
    \item \textbf{Phân bổ hàng hóa:} Trong dòng hoặc cột vừa được chọn, tìm ô có cước phí nhỏ nhất ($c_{min1}$). Phân bổ lượng hàng tối đa có thể vào ô này: $x_{ij} = \min(a_i, b_j)$.
    
    \item \textbf{Cập nhật và Gạch bỏ:} Cập nhật lại cung $a_i \leftarrow a_i - x_{ij}$ và cầu $b_j \leftarrow b_j - x_{ij}$. Trạm nào giảm về 0 trước thì gạch bỏ dòng hoặc cột đó.
    
    \item \textbf{Lặp lại:} Lặp lại từ \textbf{Bước 1} (tính lại các $\Delta$ mới cho phần bảng còn lại) cho đến khi chỉ còn lại 1 dòng hoặc 1 cột thì phân bổ nốt lượng hàng dựa vào cước phí, không cần tính $\Delta$ nữa.
\end{enumerate}

\textbf{Tổng quát hóa:} Phương pháp Vogel có thể được mô tả bằng công thức toán học như sau:
\begin{enumerate}
    \item Tính $\Delta$ cho tất cả dòng và cột chưa gạch bỏ.
    \item Chọn dòng/cột có $\max \Delta$.
    \begin{itemize}
        \begin{formula}
            \Delta_i = c_{min2} - c_{min1} \quad \text{(cho dòng $i$)}
        \end{formula}
        \begin{formula}
            \Delta_j = c_{min2} - c_{min1} \quad \text{(cho cột $j$)}
        \end{formula}
    \end{itemize}
    \item Phân bổ vào ô có cước phí nhỏ nhất ($x_{ij} = \min(p_i, t_j)$) trên dòng/cột đó.
    \item Cập nhật cung cầu và gạch bỏ nếu cần.
\end{enumerate}

Trong đó:
\begin{itemize}
    \item $c_{min1}$: Cước phí nhỏ nhất trên dòng/cột đang xét.
    \item $c_{min2}$: Cước phí nhỏ thứ hai trên dòng/cột đang xét.
    \item $\Delta$: Độ chênh lệch giữa cước phí nhỏ thứ hai và cước phí nhỏ nhất, phản ánh chi phí cơ hội nếu không chọn ô có cước phí thấp nhất.
\end{itemize}

\begin{figure}[H]
    \centering
    \resizebox{1\textwidth}{!}{%
    \begin{tikzpicture}[
        >={Stealth[scale=1.2]},
        line cap=rect,
        line join=miter,
        cell/.style={
            rectangle,
            draw=black,
            line width=0.5pt,
            minimum width=3.4cm,
            minimum height=1.6cm,
            align=center,
            anchor=center,
            inner sep=2pt
        },
        pcell/.style={
            rectangle,
            draw=none,
            minimum width=1.2cm,
            minimum height=1.6cm,
            align=center,
            anchor=center,
            font=\bfseries\color{black!70}
        },
        rcell/.style={
            rectangle,
            draw=none,
            minimum width=3.4cm,
            minimum height=0.8cm,
            align=center,
            anchor=center,
            font=\bfseries\color{black!70}
        },
        corner/.style={
            rectangle,
            draw=none,
            minimum width=3.4cm,
            minimum height=0.8cm,
            align=center,
            anchor=center,
            font=\bfseries\color{black}
        },
        ampersand replacement=\&
    ]
        % 1. Vẽ ma trận mở rộng chứa cả Delta
        \matrix (M) [
            matrix of nodes,
            nodes in empty cells,
            column sep=-0.5pt,
            row sep=-0.5pt
        ]
        {
            |[cell, alias=diag]| \& |[cell]| $\cancel{300} \rightarrow 0 : T_1$ \& |[cell]| $\cancel{500} \rightarrow \cancel{100} \rightarrow 0 : T_2$ \& |[cell]| $\cancel{100} \rightarrow 0 : T_3$ \& |[pcell]| $\Delta^{(1)}$ \& |[pcell]| $\Delta^{(2)}$ \\
            |[cell]| $\cancel{700} \rightarrow \cancel{400} \rightarrow 0: P_1$ \& |[cell, alias=cell11]| $5 \mid 300$ \& |[cell, alias=cell12]| $6 \mid 400$ \& |[cell, alias=cell13]| $4 \mid 0$ \& |[pcell]| 1 \& |[pcell]| 1 \\
            |[cell]| $\cancel{200} \rightarrow \cancel{100} \rightarrow 0 : P_2$ \& |[cell, alias=cell21]| $8 \mid 0$ \& |[cell, alias=cell22]| $9 \mid 100$ \& |[cell, alias=cell23]| $5 \mid 100$ \& |[pcell, alias=dP2_1]| 3 \& |[pcell]| 1 \\
            |[corner]| $\Delta^{(1)}$ \& |[rcell]| 3 \& |[rcell]| 3 \& |[rcell]| 1 \& \& \\
            |[corner]| $\Delta^{(2)}$ \& |[rcell, alias=dT1_2]| 3 \& |[rcell]| 3 \& |[rcell]| -- \& \& \\
        };

        % Ký hiệu x_ij ở góc dưới-trái từng ô
        \node[anchor=south west, xshift=2pt, yshift=2pt, font=\scriptsize] at (cell11.south west) {$x_{11}$};
        \node[anchor=south west, xshift=2pt, yshift=2pt, font=\scriptsize] at (cell12.south west) {$x_{12}$};
        \node[anchor=south west, xshift=2pt, yshift=2pt, font=\scriptsize] at (cell13.south west) {$x_{13}$};
        \node[anchor=south west, xshift=2pt, yshift=2pt, font=\scriptsize] at (cell21.south west) {$x_{21}$};
        \node[anchor=south west, xshift=2pt, yshift=2pt, font=\scriptsize] at (cell22.south west) {$x_{22}$};
        \node[anchor=south west, xshift=2pt, yshift=2pt, font=\scriptsize] at (cell23.south west) {$x_{23}$};

        % 2. Xử lý ô 1x1 (Đường chéo và Label)
        \draw[black!70, thick] (diag.north west) -- (diag.south east);
        \node[anchor=north east, xshift=-3pt, yshift=-3pt, font=\footnotesize, fill=white, inner sep=1pt] at (diag.north east) {Trạm thu ($T_j$)};
        \node[anchor=south west, xshift=3pt, yshift=3pt, font=\footnotesize, fill=white, inner sep=1pt] at (diag.south west) {Trạm phát ($P_i$)};

        % ---------------------------------------------------------
        % 3. KHOANH TRÒN VÀ VẼ MŨI TÊN (ĐÃ TỐI ƯU THẨM MỸ)
        % ---------------------------------------------------------
        
        % LẦN LẶP 1 (MÀU ĐỎ)
        \node[draw=red, circle, thick, inner sep=1pt, minimum size=1.8em] at (dP2_1) {};
        % Ép mũi tên đi sát mép trên của ô để không đâm qua chữ ở giữa
        \draw[->, red, dashed, thick] ([yshift=0.55cm]dP2_1.west) -- ([yshift=0.55cm]cell23.east);
        
        % LẦN LẶP 2 (MÀU XANH DƯƠNG)
        \node[draw=blue, circle, thick, inner sep=1pt, minimum size=1.8em] at (dT1_2) {};
        % Ép mũi tên đi sát mép phải của ô để không đâm qua chữ ở giữa
        \draw[->, blue, dashed, thick] ([xshift=1.4cm]dT1_2.north) -- ([xshift=1.4cm]cell11.south);

    \end{tikzpicture}
    }
    \caption{Ví dụ Phương pháp Vogel với bảng tính $\Delta$ và luồng phân bổ chi tiết}
    \label{fig:vogel-method-transport}
\end{figure}

\textbf{Diễn giải chi tiết các bước tính toán:}

\begin{enumerate}
    \item \textbf{Lần lặp 1 (\textcolor{red}{Đường đứt nét màu đỏ}):} 
    \begin{itemize}
        \item Tính $\Delta^{(1)}$: $\Delta_{P1} = 5 - 4 = 1$; $\Delta_{P2} = 8 - 5 = 3$; $\Delta_{T1} = 8 - 5 = 3$; $\Delta_{T2} = 9 - 6 = 3$; $\Delta_{T3} = 5 - 4 = 1$.
        \item $\max \Delta = 3$ (xuất hiện ở $P_2, T_1, T_2$). Chọn tùy ý dòng $P_2$ (\textcolor{red}{khoanh đỏ}).
        \item Trên dòng $P_2$, ô có cước phí nhỏ nhất là $c_{23} = 5$. Phân bổ $x_{23} = \min(100, 200) = 100$.
        \item Cập nhật: Trạm $T_3$ đã nhận đủ 100 $\rightarrow$ \textbf{Gạch bỏ cột $T_3$}. Cung $P_2$ còn lại 100.
    \end{itemize}
    
    \item \textbf{Lần lặp 2 (\textcolor{blue}{Đường đứt nét màu xanh dương}):} (Bỏ qua cột $T_3$)
    \begin{itemize}
        \item Tính $\Delta^{(2)}$: $\Delta_{P1} = 6 - 5 = 1$; $\Delta_{P2} = 9 - 8 = 1$; $\Delta_{T1} = 8 - 5 = 3$; $\Delta_{T2} = 9 - 6 = 3$.
        \item $\max \Delta = 3$ (xuất hiện ở $T_1, T_2$). Chọn cột $T_1$ (\textcolor{blue}{khoanh xanh}).
        \item Trên cột $T_1$, ô có cước phí nhỏ nhất là $c_{11} = 5$. Phân bổ $x_{11} = \min(300, 700) = 300$.
        \item Cập nhật: Trạm $T_1$ đã nhận đủ 300 $\rightarrow$ \textbf{Gạch bỏ cột $T_1$}. Cung $P_1$ còn lại 400.
    \end{itemize}

    \item \textbf{Phân bổ phần còn lại:} 
    \begin{itemize}
        \item Bảng lúc này chỉ còn duy nhất cột $T_2$. Ta không cần tính $\Delta$ nữa mà dồn lượng cung còn lại của $P_1$ và $P_2$ vào cột này.
        \item Cấp $x_{12} = \min(500, 400) = 400 \rightarrow$ \textbf{Gạch bỏ dòng $P_1$}.
        \item Cầu $T_2$ còn thiếu 100, lấy nốt từ $P_2$: $x_{22} = \min(100, 100) = 100$.
        \item Mọi cung cầu đều về 0. Thuật toán kết thúc.
    \end{itemize}
\end{enumerate}

\textbf{Kết luận:} Phương án cơ bản xuất phát thu được là:
\[ X^0 = \begin{pmatrix} 300 & 400 & 0 \\ 0 & 100 & 100 \end{pmatrix} \]
$X^0$ có 4 ô chọn = $m + n - 1 = 2 + 3 - 1 = 4 \Rightarrow X^0$ là phương án cơ bản không suy biến. \\
Tổng cước phí vận chuyển ban đầu là: 
$Z = (300 \times 5) + (400 \times 6) + (100 \times 9) + (100 \times 5) = 1500 + 2400 + 900 + 500 = 5300$.
    \section{Phương pháp thế vị đối với bài toán Min}

\begin{figure}[h!]
\centering
\begin{tikzpicture}[
    node distance=1.5cm and 2.5cm,
    box/.style={rectangle, draw, rounded corners, text width=3.5cm, align=center, minimum height=1cm, fill=blue!5},
    decision/.style={diamond, draw, aspect=2, text width=3cm, align=center, inner sep=0pt, fill=green!5},
    arrow/.style={-{Latex[length=3mm]}, thick},
    annotation/.style={text width=7cm, align=left, font=\small, draw, dashed, inner sep=10pt, fill=yellow!5}
]

% --- Khởi tạo các Node ---
\node[box] (start) {$X^0$ \\ (PACB không suy biến)};
\node[box, below=of start] (step1) {Tìm hệ thống thế vị \\ $u_i, v_j$};
\node[decision, below=of step1] (step2) {Tính hệ số ước lượng \\ $\Delta_{ij}$};
\node[box, right=2.5cm of step2] (end) {Dừng thuật toán \\ (Phương án tối ưu)};
\node[box, below=of step2] (step3) {Tìm ô vào cơ sở $(m,k)$};
\node[box, left=2.5cm of step3] (step4) {Xác định lượng hàng điều chỉnh $d$};
\node[box, above=of step4] (step5) {Tìm ô ra khỏi cơ sở $(p,q)$};

% --- Vẽ các mũi tên (Chu trình) ---
\draw[arrow] (start) -- node[right] {(6)} (step1);
\draw[arrow] (step1) -- node[right] {(1)} (step2);
\draw[arrow] (step2) -- node[right] {(2)} (step3);
\draw[arrow] (step3) -- node[above] {(3)} (step4);
\draw[arrow] (step4) -- node[left] {(4)} (step5);
\draw[arrow] (step5) |- node[near end, above] {(5)} (start);

% --- Mũi tên điều kiện dừng ---
\draw[arrow] (step2) -- node[above] {$\Delta_{ij} \le 0, \forall i,j$} (end);

% --- Khung chú thích ---
\node[annotation, right=1.5cm of start, anchor=north west] (notes) {
    \textbf{Chú thích các bước:}\\
    \textbf{(1):} Giải hệ $u_i + v_j = c_{ij}$ (cho các ô chọn).\\
    \textbf{(2):} Tính $\Delta_{ij} = u_i + v_j - c_{ij}$ (cho các ô loại).\\
    \textbf{(3):} Chọn ô vào là ô có $\max\{\Delta_{ij} > 0\} = \Delta_{mk}$.\\
    $\rightarrow$ Bổ sung lượng hàng vào ô $(m,k)$.\\
    $\rightarrow$ Lập chu trình, đánh dấu (+) và (-) xen kẽ (ô vào mang dấu +).\\
    \textbf{(4):} $d = \min \{x_{ij} \mid x_{ij} \text{ mang dấu (-) }\} = x_{pq}$.\\
    \textbf{(5):} Ô ra là ô $(p,q)$ có lượng hàng bằng 0 sau điều chỉnh.\\
    \textbf{(6):} Cập nhật phương án mới $X^1$:\\
    $x_{ij}^{mới} = \begin{cases}
    x_{ij}, & \text{nếu ô không có dấu} \\
    x_{ij} + d, & \text{nếu ô có dấu (+)} \\
    x_{ij} - d, & \text{nếu ô có dấu (-)}
    \end{cases}$
};

\end{tikzpicture}
\end{figure}